\subsection{Direct sums of cyclic modules}

Let A be a commutative ring. Recall (II, p.~220, Prop. 22) that a cyclic A-module is isomorphic to a quotient module A/a, where a is an ideal of A. We will see later (Sect. 4) that every finitely generated module over a principal ideal domain is a direct sum of finitely many cyclic modules.

PROPOSITION 1. --- Let E be a module over a commutative ring A; suppose E is a direct sum of n cyclic modules \(A/\mathfrak{a}_{k}\) ( \(1 \leqslant k \leqslant n\) ), where the \(a_{k}\) are ideals of A;

then, for each integer \(p > 0\) , the A-module \(\bigwedge^p \mathbf{E}\) is isomorphic to the direct sum of the modules \(\mathbf{A} / \mathfrak{a}_{\mathrm{H}}\) , where for each \(p\) -element subset \(\mathbf{H} = \{k_1, \dots, k_p\}\) of \([1, n]\) , the ideal \(\mathfrak{a}_{\mathrm{H}}\) is \(\sum_{j=1}^{p} \mathfrak{a}_{k_j}\) .

Let \(x_{k}\) be the canonical generator of \(A / a_{k}\) , that is the image of the unit element of A, so that E is the direct sum of the \(Ax_{i}\) ( \(1 \leqslant i \leqslant n\) ). Then we know (III, p.~515, Prop. 10) that the exterior algebra \(\bigwedge_{n} E\) is isomorphic as an A-module to the tensor product \(\otimes_{i=1}^{n} (\bigwedge(Ax_{i}))\) . Now \(\bigwedge(Ax_{i})\) is just the direct sum \(A \oplus Ax_{i}\) , since

\[
\bigwedge^ {p}
\]

the exterior product of any two elements of \(Ax_{i}\) is zero, and so \(\Lambda E\) is the direct sum of the modules \(\mathbf{M}_{\mathrm{H}} = (\mathbf{A}x_{k_{1}}) \otimes \ldots \otimes (\mathbf{A}x_{k_{p}})\) as \(H = \{k_{1}, \ldots, k_{p}\}\) runs over the set of p-element subsets of {[}1, n{]} (with \(k_{1} < \ldots < k_{p}\) ); now \(M_{H}\) is known to be isomorphic to \(A/a_{H}\) (II, p.~257, Cor. 4), which completes the proof.

We now see that, in the notation of Prop. 1, if the ideals \(\mathfrak{a}_k\) form an increasing sequence, they are completely determined by the module E. More precisely:

PROPOSITION 2. --- Let A be a commutative ring, and let E be a direct sum of n cyclic modules \(A/a_{k}\) , where the \(a_{k}\) satisfy \(a_{1} \subset a_{2} \subset \ldots \subset a_{n}\) . Then, for \(1 \leqslant p \leqslant n\) , the ideal \(a_{p}\) is the annihilator of \(\bigwedge^{p} E\) ; if \(a_{n} \neq A\) then \(\bigwedge^{p} E \neq 0\) for \(1 \leqslant p \leqslant n\) and \(\bigwedge^{m} E = 0\) for m \textgreater{} n.

Indeed, in the notation of Prop. 1, we have \(\mathfrak{a}_{\mathrm{H}} = \mathfrak{a}_{s(\mathrm{H})}\) , where \(s(\mathrm{H})\) denotes the greatest element of the subset H. Since \(s(\mathrm{H}) \geqslant p\) for every \(p\) -element subset H, and since \(s(\mathrm{H}) = p\) for \(\mathrm{H} = \{1, \dots, p\}\) , it follows that \(\mathfrak{a}_p\) is the intersection of \(\mathfrak{a}_{\mathrm{H}}\) as H varies over the set of \(p\) -element subsets of \([1, n]\) ; the ideal \(\mathfrak{a}_p\) is thus indeed the annihilator of \(\bigwedge^p E\) , by Prop. 1.

COROLLARY. --- In the notation of Prop. 2, if \(\mathfrak{a}_n \neq \mathbf{A}\) , and if E is also isomorphic to the direct sum of \(m\) cyclic modules \(\mathbf{A} / \mathfrak{a}_j'\) with \(\mathfrak{a}_1' \subset \mathfrak{a}_2' \subset \ldots \subset \mathfrak{a}_m' \neq \mathbf{A}\) , then \(m = n\) and \(\mathfrak{a}_k = \mathfrak{a}_k'\) for \(1 \leqslant k \leqslant n\) (« uniqueness of the \(\mathfrak{a}_k\) »).
% semantic-fragments: legacy-input-seam
\relax{}
